\documentclass[11pt]{article}
\usepackage[margin=1in]{geometry}
\usepackage{amsmath,amssymb,amsthm,mathtools,bm}
\usepackage{microtype}
\usepackage[T1]{fontenc}
\usepackage{lmodern}
\usepackage{booktabs,enumitem}
\usepackage[hidelinks]{hyperref}
\setlist{nosep}
\newtheorem{theorem}{Theorem}[section]
\newtheorem{corollary}[theorem]{Corollary}
\newtheorem{lemma}[theorem]{Lemma}
\newtheorem{proposition}[theorem]{Proposition}
\theoremstyle{definition}
\newtheorem{definition}[theorem]{Definition}
\newtheorem{assumption}[theorem]{Assumption}
\theoremstyle{remark}
\newtheorem{remark}[theorem]{Remark}

\newcommand{\GM}{GM}
\newcommand{\FH}{F_H}
\newcommand{\Tr}{\operatorname{Tr}}

\title{\textbf{Synopsis: Nonsingular Gravitational Collapse via Hysterical Instability}\\[0.35em]
\large Hysterical Response Theory proves a stable finite-radius remnant}
\author{Andrew Bruce Boyd}
\date{September 2026}

\begin{document}
\maketitle
\noindent\textit{Computational note.}
Proofs, exploratory experiments, and paper writing were assisted by generative AI.
Mathematical theorems stated in this paper are formally verified in Lean~4 in the
companion specifications, as faithful ports of the TeX arguments at the abstractions
recorded there, except results explicitly tagged as \emph{literature hypotheses}
(standard theorems cited from the literature and not re-proved here) or
\emph{physical inputs} (modeling assumptions outside mathematics).

\begin{abstract}
Short guide to the nonsingular-collapse paper.
HRT identifies the wake with ordinary gravity's residual response; the remnant theorem
then yields nonsingular spherical collapse.
\end{abstract}

\section{Bridge from HRT}
Hysterical Response is the residual / off-classical part of an already-existing
interaction,
$\FH=-\Tr(F\mathcal L^{-1}\mathcal S)$---for gravity, a leftover of ordinary gravity,
not a new force.
Neutrality does not kill gravitational $\FH$; closed-loop gain crossing unity makes the
channel unstable.

Classical spherical gravity cancels exterior shells.
That cancellation is organizational for the residual response.
After horizon crossing, the interior cannot reorganize the exterior residual left in the
vacated region: the wake still attracts, and relative to the collapsing surface that
attraction is \emph{outward}.

\section{Wake and criticality}
\begin{theorem}[Outward Hysterical wake]
Under gravitational HRT and horizon residual-cancellation failure, the wake
acceleration $H(R)=\hat r\cdot\mathbf g_H^{\mathrm{wake}}$ is a strictly positive
outward acceleration on the collapsing surface.
\end{theorem}

Increasing curvature drives HRT criticality, so $H$ strengthens inward and dominates
$\GM/R^{2}$ near the classical singular limit while remaining subdominant before the
transition.

\section{Remnant}
Net force $F(R)=H(R)-\GM/R^{2}$, dynamics $\ddot R=F-\gamma\dot R$.

\begin{theorem}[Stable Hysterical Remnant]
A finite $R_\ast>0$ with $H(R_\ast)=\GM/R_\ast^{2}$ exists.
If $H'(R_\ast)<-2\GM/R_\ast^{3}$ and $\gamma>0$, it is linearly asymptotically stable.
\end{theorem}

\section{Main claim}
\begin{theorem}[Nonsingular Gravitational Collapse]
Under spherical HRT, horizon residual-cancellation failure, curvature-driven
criticality, the radial reduction, restoring slope, and positive damping, spherical
collapse ends at a stable remnant $R_\ast>0$ rather than $R=0$.
In that sense, Hysterical Response Theory proves nonsingular gravitational collapse.
\end{theorem}

The load-bearing geometric input is horizon residual-cancellation failure; the wake
identification and instability engine are HRT.

\end{document}
